\documentclass[12pt, a4paper]{article}

% ── Packages ──────────────────────────────────────────────────────────
\usepackage[utf8]{inputenc}
\usepackage[T1]{fontenc}
\usepackage{amsmath, amssymb}
\usepackage{graphicx}
\usepackage{booktabs}
\usepackage{natbib}
\usepackage[margin=1in]{geometry}
\usepackage{setspace}
\usepackage{hyperref}
\usepackage{xcolor}

\doublespacing

% ── Title ─────────────────────────────────────────────────────────────
\title{Personality Traits Predict Mental Health Across Three Universities:\\
       A Multi-Study Machine Learning Investigation ($N = 1{,}559$)}

\author{
  Author A\textsuperscript{1},
  Author B\textsuperscript{1},
  Cyrus S.~H.~Ho\textsuperscript{1}
  \\[6pt]
  \textsuperscript{1}Department of Psychological Medicine,\\
  National University of Singapore
}

\date{}

\begin{document}
\maketitle

% ══════════════════════════════════════════════════════════════════════
\begin{abstract}
% ~250 words
% Background → Gap → Method → Results → Conclusion
% Key numbers: 3 studies, N=1559, 3 universities, 5 MH instruments,
%   Neuroticism #1 in 28/28 SHAP models, AUC 0.75–0.86, meta r=0.44–0.63
\end{abstract}

\textbf{Keywords:} Big Five personality, neuroticism, mental health prediction, passive sensing, SHAP, machine learning, replication

\newpage

% ══════════════════════════════════════════════════════════════════════
\section{Introduction}

\subsection{Personality and Mental Health}

The Five-Factor Model of personality \citep[FFM;][]{mccrae1992introduction, goldberg1993structure} has emerged as one of the most robust frameworks for understanding individual differences in psychological functioning. Comprising five broad trait dimensions---Neuroticism, Extraversion, Openness, Agreeableness, and Conscientiousness---the FFM provides a comprehensive taxonomy of dispositional tendencies that has been replicated across cultures, age groups, and assessment methods over more than three decades of research. Among the most consequential applications of this framework has been the study of personality--psychopathology relations, where a large and consistent body of evidence demonstrates that personality traits are powerful predictors of mental health outcomes.

The landmark meta-analysis by \citet{kotov2010linking}, synthesizing 175 studies and over 75,000 participants, established the empirical magnitude of these associations with considerable precision. Neuroticism emerged as the dominant predictor, associated with depressive disorders at $d = 1.65$ and with anxiety disorders at $d = 1.40$---effect sizes that rank among the largest in clinical psychology. The association is not limited to concurrent symptomatology. \citet{lahey2009public} argued that Neuroticism constitutes a transdiagnostic risk factor of substantial public health significance, prospectively predicting the onset of depression, anxiety, substance use disorders, and other forms of psychopathology across longitudinal studies spanning decades. Subsequent meta-analytic work confirmed that elevated Neuroticism approximately doubles the risk of developing a mental disorder even after adjusting for baseline symptoms and psychiatric history \citep{dewall2011agree}, supporting the interpretation that Neuroticism is not merely a correlate of distress but a genuine vulnerability factor.

Beyond Neuroticism, other FFM dimensions contribute meaningfully to the prediction of mental health. Conscientiousness has been identified as a protective factor, with meta-analytic evidence showing negative associations with depression ($d = -0.90$) and substance use disorders \citep{kotov2010linking}. Low Extraversion is associated with depression and social anxiety, while low Agreeableness shows links to externalizing psychopathology. \citet{roberts2007power} demonstrated that personality traits rival socioeconomic status and cognitive ability in predicting important life outcomes, including health and mortality. Theoretical accounts have conceptualized the personality--psychopathology relationship through both vulnerability models, in which traits predispose individuals to disorders, and spectrum models, in which traits and disorders represent different points on a shared continuum of functioning \citep{widiger2007plate}.

Despite this extensive literature, several important questions remain unresolved. First, the vast majority of personality--mental health research has relied on traditional statistical methods---zero-order correlations, hierarchical regression, structural equation modeling---rather than machine learning approaches that evaluate out-of-sample predictive accuracy. The distinction matters: a correlation of $r = .50$ in a training set does not guarantee that personality scores will generalize to predict symptoms in new individuals, particularly when multiple predictors are considered simultaneously and overfitting is a risk. Second, multi-study replications using identical analytic pipelines are rare. Most meta-analyses aggregate results across studies that differ in their measures, statistical methods, sample characteristics, and analytic choices, making it difficult to determine whether effect size heterogeneity reflects genuine population differences or merely methodological inconsistency. Third, clinical utility metrics---sensitivity, specificity, area under the ROC curve---are almost never reported, leaving unclear whether the statistical associations between personality and mental health translate into practically useful screening performance.

\subsection{Passive Behavioral Sensing}

While the personality--psychopathology literature has accumulated evidence over decades, a newer research tradition has sought to predict mental health using passively collected behavioral data from smartphones and wearable devices. \citet{insel2018digital} characterized the smartphone as ``a potential game changer for psychiatry,'' envisioning continuous, passive collection of behavioral signals---movement patterns, sleep rhythms, social interactions, screen use---as a pathway to objective, real-time psychiatric assessment. This vision of \textit{digital phenotyping} has been elaborated in comprehensive reviews \citep{torous2021growing, onnela2016harnessing, mohr2017personal} and has attracted substantial investment in dedicated research centers, government-funded initiatives, and commercial platforms.

Empirical support for the approach emerged from early studies such as \citet{saeb2015mobile}, who reported correlations between GPS-derived mobility features and PHQ-9 depression scores, and the landmark StudentLife project \citep{wang2014studentlife}, which demonstrated that smartphone sensing could track behavioral changes across a college academic term. Subsequent work linked smartphone and wearable data to depression \citep{chikersal2021detecting, wang2018sensing}, anxiety \citep{sano2018identifying}, and general well-being \citep{ben2015next}. \citet{harari2016using} outlined how smartphones could serve as ``pocket laboratories'' for psychological science, and \citet{jacobson2019digital} proposed that digital biomarkers could supplement or even supplant traditional symptom measures.

An implicit assumption underlying much of this work is that continuous behavioral data should provide information \textit{beyond} what traditional questionnaires capture. Smartphones record what people \textit{do}---their mobility, sleep, social interactions, digital habits---rather than what they report on subjective scales. The data are collected passively and continuously, requiring no ongoing participant effort. These features have made behavioral sensing an attractive complement---or potential replacement---for conventional self-report assessment. Yet the question of whether sensing data genuinely add predictive value over and above brief, well-validated questionnaires has received surprisingly little direct attention.

\subsection{Gaps and Present Study}

Despite the maturity of the personality--mental health literature and the rapid growth of digital phenotyping, several critical gaps remain at the intersection of these two traditions.

\textit{Gap 1: Multi-study replication with identical ML pipelines.} Existing machine learning studies of personality and mental health are typically conducted in single samples, using study-specific analytic choices. The absence of multi-study replications with standardized pipelines makes it difficult to assess the robustness and generalizability of reported effects. Meta-analyses partially address this concern but cannot control for analytic heterogeneity across primary studies.

\textit{Gap 2: Clinical utility evaluation.} Although the statistical association between personality traits and mental health is well established, translation to clinical screening contexts---characterized by AUC, sensitivity, and specificity at established cutoffs---remains almost entirely absent from the literature. Without these metrics, it is unclear whether personality-based models can achieve performance levels useful for population-level mental health screening.

\textit{Gap 3: Incremental validity of behavioral sensing over personality.} Studies in the digital phenotyping literature typically report sensing--outcome associations without including personality as a baseline predictor \citep{chikersal2021detecting, wang2018sensing, saeb2015mobile}, while personality--psychopathology studies rarely incorporate passively sensed features \citep{kotov2010linking, lahey2009public}. As a result, the incremental validity of behavioral sensing beyond personality---tested with appropriate multiple comparison correction---remains essentially unknown. Growing concerns about the robustness of sensing-based predictions reinforce this gap: \citet{muller2021depression} found that GPS-based depression models dropped from AUC = 0.82 to 0.57 in larger samples, and the GLOBEM benchmark \citep{xu2022globem} reported near-chance cross-dataset generalization. \citet{dasswain2022semantic} attributed such failures to a fundamental \textit{semantic gap} between low-level sensor signals and psychological constructs.

\textit{Gap 4: SHAP versus traditional methods.} Interpretable machine learning techniques such as SHAP \citep{lundberg2017unified} have been applied to mental health prediction models, but it remains unclear whether data-driven feature importance rankings discover substantively different patterns than traditional regression coefficients. If SHAP rankings converge with standardized betas, this would suggest that the personality--mental health relationship is predominantly linear and that complex interpretability methods provide communicative rather than substantive gains.

The present study addresses these four gaps through a multi-study machine learning investigation of personality traits as predictors of mental health outcomes. We analyze three publicly available longitudinal datasets---StudentLife \citep{wang2014studentlife}, NetHealth, and GLOBEM \citep{xu2022globem}---totaling $N = 1{,}559$ undergraduates across three U.S.\ universities (Dartmouth College, University of Notre Dame, University of Washington). All three datasets contain both Big Five personality assessments and passive behavioral sensing data (smartphone and wearable), enabling direct comparison of these two measurement modalities within the same participants. A unified machine learning pipeline (Ridge regression, Elastic Net, Random Forest, SVR, and logistic regression with Optuna hyperparameter tuning, repeated cross-validation, and permutation testing) is applied identically across all datasets.

We organize our investigation around four aims:

\begin{enumerate}
    \item[\textbf{Aim 1.}] Quantify the predictive accuracy of Big Five personality traits for mental health outcomes (depression, anxiety, stress, loneliness) using machine learning with out-of-sample evaluation, replicated across three independent studies.
    \item[\textbf{Aim 2.}] Evaluate clinical utility by reporting AUC, sensitivity, and specificity at established clinical cutoffs for personality-based screening models.
    \item[\textbf{Aim 3.}] Test the incremental validity of passive behavioral sensing beyond personality using nested model comparisons with false discovery rate correction.
    \item[\textbf{Aim 4.}] Compare SHAP-derived feature importance rankings against traditional regression coefficients to assess whether machine learning interpretability reveals structure beyond what linear methods capture.
\end{enumerate}

By applying an identical pipeline across three datasets, we provide a multi-study replication that moves beyond single-sample findings to assess the robustness and boundaries of personality-based mental health prediction. We further examine whether more complex models (multilayer perceptrons) improve upon simple penalized regression, whether demographic covariates attenuate personality effects, and whether the pattern of results holds across different mental health constructs and personality instruments (BFI-44 vs.\ BFI-10). The results speak to the practical question of what the simplest, most cost-effective approach to mental health screening might be---and whether the field's investment in continuous behavioral sensing provides meaningful returns beyond what a brief personality questionnaire already delivers.

% ══════════════════════════════════════════════════════════════════════
\section{Method}

\subsection{Datasets}

We conducted secondary analyses of three publicly available longitudinal datasets from United States universities, comprising a total of $N = 1{,}559$ undergraduate participants. Table~1 summarizes the key characteristics of each study. All three studies received institutional ethical approval from their respective universities, and we refer readers to the original publications for detailed ethical review procedures.

\subsubsection{Study 1: StudentLife (Dartmouth College)}

The StudentLife dataset \citep{wang2014studentlife} comprises data from a 10-week smartphone sensing study conducted at Dartmouth College during the Fall 2013 academic term. Of 48 enrolled participants, $N = 28$ had complete Big Five personality and cumulative GPA records and formed our analytic sample. Participants were undergraduate and graduate students who carried Android smartphones instrumented with continuous sensing software throughout the term. Six outcome measures were available: cumulative GPA, depressive symptoms (PHQ-9), perceived stress (PSS), loneliness, flourishing, and negative affect (PANAS-NA).

\subsubsection{Study 2: NetHealth (University of Notre Dame)}

The NetHealth study is a large-scale longitudinal investigation of college student health conducted at the University of Notre Dame from 2015 to 2019. Participants were first-year students who wore Fitbit wearable devices and completed periodic surveys over multiple academic years. The full sample comprised $N = 722$ participants: 220 with both BFI-44 and GPA records, and 498 with both BFI-44 and at least one mental health instrument. Of these, $N = 365$ had complete behavioral sensing data (Fitbit activity, sleep, and communication logs). Demographic information was available for this study, including gender, native status, parental education, and family income. Mental health outcomes included the CES-D (depression), STAI-Trait (trait anxiety), and BAI (Beck anxiety).

\subsubsection{Study 3: GLOBEM (University of Washington)}

The GLOBEM dataset \citep{xu2022globem} is a multi-year longitudinal study conducted at the University of Washington from 2018 to 2021, spanning four annual cohorts (INS-W\_1 through INS-W\_4). The total sample comprised $N = 809$ participants (800 with BFI-10 personality data; 786 with BDI-II depression scores). Participants wore Fitbit devices and carried smartphones instrumented with sensing software. The third cohort (INS-W\_3, 2020--2021) overlapped with the COVID-19 pandemic, providing a natural test of robustness to environmental disruption. Each cohort included approximately 200 participants. Mental health outcomes included the BDI-II (depression), STAI-State (state anxiety), PSS-10 (perceived stress), CESD-10 (depression), and UCLA Loneliness Scale. Weekly PHQ-4 assessments were available for a subset ($N = 540$) over approximately 10 weeks, enabling trajectory analyses. No GPA data were available in this dataset.

\subsection{Measures}

\subsubsection{Big Five Personality}

Personality was assessed using the Big Five Inventory. Studies~1 and~2 used the 44-item BFI \citep[BFI-44;][]{john1999big}, which measures Openness, Conscientiousness, Extraversion, Agreeableness, and Neuroticism on 5-point Likert scales. Internal consistency (Cronbach's $\alpha$) in the present samples ranged from .67 to .86 (Study~1) and .69 to .87 (Study~2). Study~3 used the 10-item short form \citep[BFI-10;][]{rammstedt2007measuring}, which contains two items per dimension. Because the BFI-10 yields lower reliability ($\alpha \approx .65$), we conducted disattenuation analyses to assess the impact of measurement error on observed effect sizes (see Supplementary Materials).

\subsubsection{Mental Health Outcomes}

Multiple validated instruments assessed mental health across the three studies, providing convergent measurement of depression, anxiety, and stress constructs:

\begin{itemize}
    \item \textbf{Depression.} The Patient Health Questionnaire-9 \citep[PHQ-9;][]{kroenke2001phq9} was used in Study~1 (9~items; clinical cutoff $\geq 10$). The Center for Epidemiological Studies Depression Scale \citep[CES-D;][]{radloff1977cesd} was used in Study~2 (20~items; cutoff $\geq 16$). The Beck Depression Inventory--II \citep[BDI-II;][]{beck1996bdi} was used in Study~3 (21~items; mild cutoff $\geq 14$, moderate cutoff $\geq 20$). The 10-item CES-D (CESD-10) was also available in Study~3.
    \item \textbf{Anxiety.} The State-Trait Anxiety Inventory \citep[STAI;][]{spielberger1983manual} was used in Studies~2 (Trait subscale; cutoff $\geq 45$) and~3 (State subscale). The Beck Anxiety Inventory \citep[BAI;][]{beck1988inventory} was used in Study~2.
    \item \textbf{Stress.} The Perceived Stress Scale \citep[PSS;][]{cohen1983global} was used in Study~1 (14~items). The PSS-10 (10-item version) was used in Study~3 (cutoff $\geq 20$).
    \item \textbf{Other outcomes.} The UCLA Loneliness Scale \citep{russell1996ucla} was available in Study~3. The PHQ-4 \citep{lowe2010phq4} provided weekly depression and anxiety screening in Study~3. Cumulative GPA from registrar records was available in Studies~1 and~2.
\end{itemize}

\subsubsection{Behavioral Features}

Passively sensed behavioral data were extracted from smartphones and wearable devices. The number and type of sensing modalities varied across studies, reflecting differences in the original study designs:

\textit{Study~1} collected data from 13 smartphone sensing modalities, including accelerometer-derived activity (stationary, walking, running), microphone-based audio classification (silence, voice, noise), conversation detection (frequency, duration), GPS location (entropy, time at home, distance traveled), screen usage (unlock frequency, duration), Bluetooth proximity (unique devices), Wi-Fi (unique access points), communication logs (calls, SMS), and app usage. This yielded 87 raw behavioral features, which were reduced to 8 orthogonal principal component analysis (PCA) composites capturing Mobility, Digital Usage, Audio Environment, Social Activity, Communication, Sleep, Screen Habits, and App Usage patterns.

\textit{Study~2} collected data from Fitbit wearables and communication logs, yielding 28 raw features: 11 activity features (daily steps, distance, calories, active minutes by intensity, sedentary minutes, step variability), 8 sleep features (total sleep time, time in bed, efficiency, awakenings, onset latency, variability, wake after sleep onset, regularity), 7 communication features (incoming/outgoing calls, call duration, SMS counts, unique contacts), and 2 quality indicators. These were reduced to 3 PCA composites: Activity-Sleep, Communication, and Quality.

\textit{Study~3} collected data from Fitbit wearables, smartphones, and GPS, yielding 19 main features: 4 activity features (daily steps, active minutes, sedentary minutes, step variability), 4 sleep features (total sleep time, efficiency, onset time, wake time), 4 call features (incoming/outgoing counts, duration, unique contacts), 3 screen features (on-frequency, duration, first/last unlock), and 4 location features (unique locations, entropy, time at home, distance traveled). These were reduced to 5 PCA composites: Activity, Sleep, Communication, Screen, and Location. As a sensitivity analysis, we also extracted 2,597 features using the RAPIDS pipeline, filtered to 1,258 after removing zero-variance columns.

PCA was applied within each study to address multicollinearity and the high feature-to-sample ratio. The number of retained components was determined by cumulative variance explained, and all PCA transformations were computed within cross-validation folds to prevent data leakage.

\subsection{Analytic Strategy}

All analyses were implemented in Python 3 using scikit-learn \citep{pedregosa2011scikit}. The same analytic pipeline was applied uniformly across all three studies to facilitate direct comparison.

\subsubsection{Machine Learning Pipeline}

We evaluated four regression algorithms spanning linear and non-linear approaches: Ridge regression \citep{hoerl1970ridge}, Elastic Net \citep{zou2005regularization}, Random Forest \citep{breiman2001random}, and Support Vector Regression \citep[SVR;][]{drucker1997svr}. For each outcome, three feature sets were compared: personality traits only, behavioral composites only, and personality plus behavior combined.

Hyperparameters were tuned using Optuna Bayesian optimization \citep{akiba2019optuna} with a Tree-structured Parzen Estimator (TPE) sampler and 30 optimization trials. Cross-validation followed a repeated $K$-fold scheme: leave-one-out cross-validation (LOO-CV) for Study~1 ($N = 28$, where standard $K$-fold would yield unstable variance estimates) and $5$-fold $\times$ $5$-repeat cross-validation for Studies~2 and~3. Within each fold, features were standardized using \texttt{StandardScaler} fit exclusively on training data to prevent information leakage. For supplementary analyses, Ridge regression with a fixed regularization parameter ($\alpha = 1.0$) was used to ensure consistency across comparisons; for main analyses, \texttt{RidgeCV} with built-in cross-validated $\alpha$ selection was employed.

Model performance was assessed using out-of-fold $R^2$ (coefficient of determination). Statistical significance was evaluated via permutation tests (1,000 permutations for Studies~2 and~3; 199 permutations for Study~1 given computational constraints with LOO-CV), in which outcome labels were randomly shuffled and the full modeling pipeline re-executed to construct a null distribution of $R^2$ values.

\subsubsection{SHAP Interpretability}

To interpret model predictions, we computed Shapley Additive Explanations \citep[SHAP;][]{lundberg2017unified} for every trained model. TreeSHAP (exact computation) was used for tree-based models (Random Forest), and KernelSHAP (model-agnostic approximation) was used for linear and kernel-based models (Ridge, Elastic Net, SVR). For each outcome, mean absolute SHAP values were computed across all test-fold predictions to rank feature importance. Cross-model consistency was quantified using Kendall's $\tau$ rank correlation between SHAP importance rankings across all model pairs. To validate SHAP against traditional methods, we compared SHAP-derived rankings with zero-order correlations and standardized OLS regression coefficients.

\subsubsection{Clinical Classification}

To evaluate clinical utility, we dichotomized mental health outcomes at established clinical cutoffs (CES-D $\geq 16$, STAI $\geq 45$, BDI-II $\geq 14$ and $\geq 20$, PSS $\geq 20$) and trained binary classifiers using Logistic Regression and Random Forest. Classification followed a repeated stratified $10$-fold $\times$ $10$-repeat cross-validation scheme to ensure stable estimates with preserved class proportions. Discrimination was assessed via area under the receiver operating characteristic curve (AUC) with 95\% percentile bootstrap confidence intervals. Optimal classification thresholds were selected using Youden's $J$ index (sensitivity + specificity $- 1$), and sensitivity at fixed specificity of 0.80 was reported for clinical relevance.

\subsubsection{Incremental Validity}

To test whether behavioral sensing adds predictive value beyond personality, we conducted nested $F$-tests comparing personality-only models against personality-plus-behavior models within an OLS regression framework. The null hypothesis was that the additional behavioral predictors do not explain variance beyond personality traits. Because multiple comparisons were conducted across outcomes and studies, $p$-values were corrected using the Benjamini--Hochberg false discovery rate \citep[BH-FDR;][]{benjamini1995controlling} procedure across all tests.

\subsubsection{Meta-Analysis}

To synthesize effect sizes across studies, we conducted random-effects meta-analyses using the DerSimonian--Laird estimator \citep{dersimonian1986meta}. Pearson correlations between personality traits and outcomes were transformed to Fisher's $z$ prior to pooling, and pooled estimates were back-transformed for interpretation. Heterogeneity was assessed using the $I^2$ statistic, where values of 25\%, 50\%, and 75\% indicate low, moderate, and high heterogeneity, respectively \citep{cohen1988statistical}. Forest plots display study-level and pooled effect sizes with 95\% confidence intervals.

\subsubsection{Demographic Controls}

In Study~2, where demographic data were available, we conducted hierarchical regression analyses to assess whether personality effects survive after controlling for demographic confounders. The regression was structured in three steps: (1)~demographics only (gender, native status, parental education, family income), (2)~demographics plus personality traits, and (3)~demographics plus personality plus behavioral composites. The incremental $\Delta R^2$ at each step was tested for significance. Study~3 did not collect demographic data beyond cohort membership, and Study~1 was too small ($N = 28$) to support covariate-adjusted models.

% ══════════════════════════════════════════════════════════════════════
\section{Results}

\subsection{Head-to-Head: Personality Questionnaires vs Passive Sensing}

Table~2 presents the primary comparison: out-of-fold $R^2$ with 95\% bootstrap confidence intervals for personality-only versus sensing-only models across all 15 study--outcome combinations in three datasets. Personality questionnaires outperformed passive behavioral sensing in 14 of 15 comparisons (93\%). The median advantage was $\Delta R^2 = .21$, with personality-only $R^2$ values ranging from near zero (S2 Loneliness: $R^2 = .01$) to substantial (S2 STAI-Trait: $R^2 = .52$ [95\% CI: .41, .62]). In contrast, sensing-only models produced $R^2 \leq 0$ in 12 of 15 comparisons, indicating worse-than-chance prediction.

The strongest personality effects emerged for anxiety-related constructs. In Study~2, personality explained 52.2\% of variance in trait anxiety ($R^2 = .52$ [.41, .62]) and 27.9\% in depression ($R^2 = .28$ [.09, .39]). In Study~3, personality explained 20.0\% of variance in state anxiety ($R^2 = .20$ [.08, .26]) and 14.3\% in perceived stress ($R^2 = .14$ [${-.00}$, .22]). Even in the smallest study ($N = 27$), personality predicted perceived stress ($R^2 = .39$ [.01, .72]) and loneliness ($R^2 = .20$ [${-.19}$, .50]).

Behavioral sensing produced only one positive $R^2$ value across all studies: S1 PHQ-9 ($R^2 = .20$ [${-.65}$, .82]), the sole comparison in which sensing outperformed personality. This was also the only outcome for which personality yielded a negative $R^2$ ($-.55$), likely attributable to overfitting in the small sample ($N = 27$) with leave-one-out cross-validation. Notably, the confidence intervals for both estimates in this comparison were extremely wide, spanning from negative to positive values, underscoring the instability of estimates at this sample size.

For academic performance, personality predicted GPA in Study~2 ($R^2 = .002$ [${-.32}$, .10]) and Study~1 ($R^2 = -.17$), though neither estimate was statistically distinguishable from zero. Sensing performed even worse for GPA, with $R^2 = -1.85$ in Study~1 and $R^2 = -.04$ in Study~2.

\subsection{Meta-Analytic Synthesis}

Table~3 and Figure~1 present random-effects meta-analytic estimates pooling personality--outcome correlations across studies. Sensing features were not meta-analyzed because $R^2 \leq 0$ in most study--outcome combinations, precluding meaningful pooled estimates.

The strongest meta-analytic effect was for Neuroticism predicting anxiety and stress (pooled $r = .63$ [95\% CI: .38, .80], $k = 3$, $N = 1{,}324$, $p < .001$), followed by Neuroticism predicting depression (pooled $r = .44$ [.24, .61], $k = 3$, $N = 1{,}304$, $p < .001$). Conscientiousness showed a moderate negative association with depression (pooled $r = -.21$ [${-.29}$, ${-.12}$], $k = 3$, $N = 1{,}304$, $p < .001$) with comparatively low heterogeneity ($I^2 = 44.8\%$, $p = .163$). Conscientiousness also predicted GPA (pooled $r = .38$ [.06, .62], $k = 2$, $N = 248$, $p = .019$), demonstrating trait-specific dissociation: Neuroticism was the primary predictor of mental health, whereas Conscientiousness was the primary predictor of academic performance.

Heterogeneity was high for most Neuroticism associations: $I^2 = 91.8\%$ for Neuroticism--Depression ($p < .001$), $I^2 = 96.2\%$ for Neuroticism--Anxiety ($p < .001$), and $I^2 = 94.4\%$ for Neuroticism--Loneliness ($p < .001$). The Neuroticism--Loneliness association was not significant (pooled $r = .13$ [${-.15}$, .40], $p = .348$), reflecting inconsistent directions across studies (S1: $r = .50$; S2: $r = -.14$; S3: $r = .17$). Extraversion--Loneliness was similarly non-significant and heterogeneous (pooled $r = -.14$ [${-.46}$, .21], $I^2 = 96.6\%$). The Neuroticism--Perceived Stress association was large but imprecise (pooled $r = .58$ [.07, .85], $k = 2$, $I^2 = 88.1\%$), reflecting the combination of a very strong effect in Study~1 ($r = .76$, $N = 27$) with a more moderate effect in Study~3 ($r = .37$, $N = 800$).

The high heterogeneity in most associations likely reflects genuine differences across studies in sample characteristics, personality measure length (BFI-44 vs.\ BFI-10), and outcome instruments, rather than sampling error alone. Despite this heterogeneity, the direction of effects was consistent for Neuroticism--Depression and Neuroticism--Anxiety across all three studies.

\subsection{Clinical Binary Classification}

Table~4 presents classification performance for detecting clinically elevated symptom levels. We evaluated personality-only, personality-plus-behavior, and behavior-only models using logistic regression and random forest classifiers with repeated stratified cross-validation.

In Study~2, where both personality and behavioral data were available for the same participants, personality-only models achieved AUC = .75 [95\% CI: .62, .88] for CES-D $\geq$ 16 and AUC = .79 [.68, .90] for STAI $\geq$ 45. Adding behavioral features improved discrimination modestly: AUC = .83 [.70, .93] for CES-D $\geq$ 16 and AUC = .86 [.71, .96] for STAI $\geq$ 45. However, behavior-only models performed near chance (AUC = .55 [.40, .70] and AUC = .58 [.42, .75]), confirming that personality carried the predictive signal.

In Study~3, classification performance was more modest. Personality-only models achieved AUC = .65 [.55, .76] for BDI-II $\geq$ 14, AUC = .66 [.55, .79] for BDI-II $\geq$ 20, and AUC = .69 [.58, .77] for PSS $\geq$ 20. Combined personality-plus-behavior models showed marginal improvement (AUC = .67 [.52, .78] for BDI-II $\geq$ 14, AUC = .69 [.52, .82] for BDI-II $\geq$ 20, AUC = .70 [.59, .81] for PSS $\geq$ 20). Behavior-only models again performed near chance (AUC = .57--$.60$).

DeLong tests comparing personality-only versus personality-plus-behavior AUC were conducted for all five classification targets. None survived Benjamini--Hochberg FDR correction (0/5; all $p_{\text{FDR}} > .44$), indicating that the observed improvements from adding behavioral sensing were not statistically reliable.

Sensitivity and specificity patterns were consistent across studies. At optimal Youden thresholds, personality-only models achieved high sensitivity (range: .71--$.85$) at the cost of moderate specificity (range: .60--$.70$). In Study~2, the combined model for STAI $\geq$ 45 showed the best balance: sensitivity = .85 [.56, 1.00] and specificity = .80 [.57, 1.00], with PPV = .63 and NPV = .95, suggesting potential utility as a screening instrument in this specific context.

\subsection{Incremental Validity of Behavioral Sensing}

Table~5 reports nested $F$-tests evaluating whether behavioral sensing composites add significant variance beyond personality traits in OLS regression. Across eight study--outcome combinations in Studies~2 and~3 (Study~1 was excluded due to insufficient sample size), only one survived Benjamini--Hochberg FDR correction: S3 BDI-II ($\Delta R^2 = .024$, $F(5, 577) = 3.31$, $p_{\text{FDR}} = .047$). The remaining seven tests showed non-significant increments, with $\Delta R^2$ values ranging from .001 (S2 BAI) to .018 (S3 STAI-State and S3 CESD), all $p_{\text{FDR}} > .08$.

Even in the one significant case, the incremental variance was small in absolute terms ($\Delta R^2 = .024$), representing a 2.4 percentage-point increase from personality-only $R^2 = .130$ to full-model $R^2 = .154$. The overall pattern was clear: behavioral sensing added negligible predictive information beyond what five self-report personality items already captured.

In Study~2, behavioral sensing contributed $\Delta R^2 = .009$ for CES-D, $\Delta R^2 = .002$ for STAI-Trait, and $\Delta R^2 = .003$ for BAI, none approaching significance after FDR correction. This was notable given that Study~2 had the richest behavioral feature set (Fitbit activity, sleep, and communication data) and the largest combined sample ($N = 363$ with complete behavioral data).

\subsection{Methodological Checks}

Several supplementary analyses assessed the robustness of the primary findings (see Supplementary Materials for full results of all 41 analyses).

\subsubsection{SHAP and OLS Agreement}

SHAP-derived feature importance rankings closely matched standardized OLS regression coefficients. Across seven personality--outcome combinations in Studies~2 and~3, Kendall's $\tau$ between $\beta$ and SHAP rankings ranged from .80 to 1.00 (mean $\tau = .94$). Top-ranked feature agreement was perfect: Neuroticism was identified as the strongest predictor by both methods in 6 of 7 comparisons, with Extraversion ranked first for UCLA Loneliness. This convergence indicates that the personality--mental health relationship is predominantly linear, and that SHAP values provide communicative rather than substantively novel information over traditional regression weights.

\subsubsection{Neural Network Comparison}

A two-layer multilayer perceptron (MLP) with Optuna-tuned hyperparameters consistently underperformed Ridge regression across all outcomes (Table~S1). In Study~1, MLP showed catastrophic overfitting ($R^2 = -4.06$ for Flourishing; $R^2 = -.47$ for Loneliness). In Studies~2 and~3, MLP $R^2$ values were systematically lower than Ridge: for example, S2 STAI-Trait MLP $R^2 = .49$ versus Ridge $R^2 = .52$; S3 STAI-State MLP $R^2 = .14$ versus Ridge $R^2 = .20$. This pattern reinforces the conclusion that the personality--mental health association is linear, and that additional model complexity introduces noise without capturing meaningful non-linear structure.

\subsubsection{Demographic Controls}

In Study~2, where demographic covariates (gender, native status, parental education, family income) were available, personality traits continued to explain substantial incremental variance beyond demographics: $\Delta R^2 = .33$ for CES-D, $\Delta R^2 = .52$ for STAI-Trait, and $\Delta R^2 = .20$ for BAI (all $p_{\text{FDR}} < .001$). In contrast, behavioral sensing added $\Delta R^2 < .01$ beyond demographics and personality for all three outcomes (all $p_{\text{FDR}} > .24$). These results confirm that personality effects are not reducible to demographic confounds and that sensing provides no incremental benefit even after accounting for shared demographic variance.

\subsubsection{Cross-Study Consistency}

Conscientiousness was identified as the top SHAP predictor of GPA across all four model types (Elastic Net, Ridge, Random Forest, SVR) in both Studies~1 and~2, demonstrating trait-specific consistency: Neuroticism dominated mental health prediction while Conscientiousness dominated academic prediction. Furthermore, Neuroticism was ranked as the top predictor for all mental health outcomes across all model types in Studies~2 and~3 (8/8 study--construct--model combinations), confirming the robustness of this finding to analytic choices.

% ══════════════════════════════════════════════════════════════════════
\section{Discussion}

The present study provides, to our knowledge, the first systematic head-to-head comparison of personality questionnaires and passive behavioral sensing for mental health prediction, replicated across three independent university samples totaling $N = 1{,}559$ participants. The results offer a clear answer to RQ1: a brief personality questionnaire consistently and substantially outperforms weeks of continuous passive sensing. They also address RQ2 and RQ3: sensing adds negligible incremental validity (1/8 comparisons significant after FDR correction), and more complex machine learning architectures cannot compensate for the fundamental signal limitation in behavioral features. We discuss these findings in five sections, moving from the core empirical result through its nuances, implications, and boundaries.

\subsection{An Asymmetric Contest: Questionnaires vs.\ Sensing}

The magnitude of the personality--sensing disparity was striking. Personality questionnaires explained meaningful variance in mental health outcomes across all three studies ($R^2 = .09$--.52), whereas sensing-only models produced $R^2 \leq 0$ in 12 of 15 comparisons---worse than predicting the sample mean. The sole exception (S1 PHQ-9, $R^2 = .20$) occurred in a sample of 27 participants, produced confidence intervals spanning from $-.65$ to $.82$, and did not replicate in the larger studies. In clinical classification, personality-only models achieved AUC = .65--.79, whereas behavior-only models hovered near chance (AUC = .53--.60). Two self-report items---``Is depressed, blue'' and ``Is emotionally stable, not easily upset''---captured more between-person mental health variance ($R^2 = .36$ for CES-D) than 28 sensing features aggregated over weeks ($R^2 = -.16$). The cost-effectiveness asymmetry is difficult to overstate: 10 seconds of self-report versus weeks of data collection and \$100+ in device costs.

This disparity is not, we argue, primarily a technology problem. It reflects a construct mismatch between what sensors measure and what mental health screening requires. Personality questionnaires directly probe stable dispositional tendencies---precisely the trait-level constructs (Neuroticism, emotional stability) that are theoretically and empirically linked to psychopathology \citep{kotov2010linking, lahey2009public}. Passive sensing, by contrast, captures moment-to-moment behavioral manifestations (step counts, screen unlocks, sleep duration) that are influenced by a host of situational factors---weather, academic schedules, social obligations---in addition to any underlying trait signal. \citet{dasswain2022semantic} formalized this intuition as the ``semantic gap'' between low-level sensor signals and the higher-order psychological constructs they are intended to index, arguing that increasing data volume cannot close a fundamentally inferential chasm. Our results across 1,559 participants and 15 outcome comparisons provide strong empirical support for this theoretical account. Even \citet{stachl2020predicting}, using rich smartphone data from 743 participants to predict personality itself, achieved a maximum correlation of $r = .37$---and predicting mental health from behavior requires traversing an additional inferential step from behavioral proxies to personality to psychopathology.

\subsection{When Does Sensing Help?}

Our findings do not support a blanket dismissal of passive sensing, but they sharply constrain the conditions under which it contributes. In Study~2, the combined personality-plus-behavior model achieved AUC = .83--.86 for depression and trait anxiety classification, a moderate improvement over personality alone (AUC = .75--.79). Critically, this gain was driven by communication metadata (call and SMS patterns) available in the NetHealth dataset---features that capture social interaction frequency and reciprocity, which are more semantically proximal to mental health constructs than step counts or sleep duration. Study~3, which lacked communication features, showed no comparable benefit from adding behavioral sensing.

At the individual level, our idiographic analyses revealed that 17\% of participants showed person-specific sensing--mood associations with $R^2 > .30$, and per-person EMA correlations averaged $|r| = .35$. These findings suggest that the sensing--mood link is real but heterogeneous: it is not absent for everyone, but it varies in magnitude and even direction across individuals, precluding a single population-level model from capturing the signal. This distinction between \textit{nomothetic} screening (one model for all) and \textit{idiographic} monitoring (personalized models) is consequential. Sensing may have genuine value for within-person state tracking---detecting when an individual deviates from their own behavioral baseline---even as it fails for between-person trait screening. Our dose-response analysis further clarified the boundary: sensing $R^2$ remained flat from 7 to 92 days of data collection, confirming that the between-person signal is not merely buried in short measurement windows but is fundamentally absent at the population level.

\subsection{Implications for Digital Phenotyping Research}

These results carry several implications for how the digital phenotyping field designs and evaluates studies. First, we propose that a \textbf{personality questionnaire baseline} should become standard practice in sensing research. Most published sensing studies report prediction accuracy without including a simple questionnaire comparator \citep{chikersal2021detecting, wang2018sensing, saeb2015mobile}, which makes it impossible to evaluate whether sensing provides unique information or merely recapitulates variance already accessible through a five-minute self-report. Reporting a sensing AUC of .65 in isolation invites a favorable interpretation; the same number compared against a personality-only AUC of .75 in the same sample tells a very different story.

Second, the scale of our study contextualizes the broader literature. With $N = 1{,}559$, our combined sample exceeds the median sample size in the digital phenotyping field by an order of magnitude. \citet{busshart2026scoping}, in a scoping review of digital phenotyping studies for depression, reported a median sample size of approximately 60 participants. Small samples amplify noise and generate optimistic performance estimates that fail to replicate---a pattern exemplified by our own Study~1 ($N = 27$), where the sole positive sensing result did not survive replication in larger datasets. Converging evidence supports this concern: \citet{muller2021depression} demonstrated that GPS-based depression prediction dropped from AUC = .82 to .57 when moving from small to large samples, the GLOBEM benchmark \citep{xu2022globem} reported near-chance cross-cohort generalization across 497 participants, and \citet{adler2022machine} found that passive sensing models trained on one longitudinal study failed to generalize to others. Our findings add a further dimension to this sobering picture: even within each dataset, with models trained and evaluated on the same participants via cross-validation, sensing features contain insufficient between-person signal to outperform a brief questionnaire.

Third, we recommend that every study claiming incremental value for sensing features should apply \textbf{false discovery rate correction} when testing multiple outcomes. In our data, 3 of 8 incremental validity tests were nominally significant ($p < .05$), but only 1 survived Benjamini--Hochberg correction---a reminder that uncorrected $p$-values across multiple comparisons inflate the apparent contribution of behavioral sensing.

\subsection{Simple Models Suffice}

A secondary but methodologically important finding is that model complexity provided no benefit. SHAP-derived feature importance rankings were virtually identical to standardized OLS regression coefficients (mean Kendall's $\tau = .94$ across seven personality--outcome comparisons), indicating that the penalized linear models captured the same structure that ordinary regression reveals. The multilayer perceptron consistently underperformed Ridge regression across all outcomes, often catastrophically so in small samples (S1 MLP $R^2 = -4.06$ for Flourishing). Optuna hyperparameter tuning, applied systematically across all model types, did not uncover parameter configurations that materially improved performance.

These results suggest that the personality--mental health relationship is predominantly linear, and that the field's investment in increasingly complex architectures---deep learning, transformers, attention mechanisms applied to EMA sequences---may not address the fundamental constraint. When the underlying signal is linear and low-dimensional (five personality traits predicting mental health), methodological complexity adds estimation noise without capturing meaningful non-linear structure. The appropriate level of analytic sophistication should match the complexity of the data-generating process, not the ambitions of the modeler.

\subsection{Limitations}

Several limitations constrain the generalizability of our findings. First, all three samples comprise undergraduate students at selective U.S.\ universities (Dartmouth, Notre Dame, University of Washington), restricting inference to relatively affluent, young, and predominantly White populations. Whether the personality--sensing comparison generalizes to clinical populations, older adults, or non-Western samples remains an open question.

Second, our design is cross-sectional with respect to the core prediction models. Although supplementary longitudinal and trajectory analyses confirmed that personality predicts mean symptom levels but not within-person change over time---consistent with a trait interpretation---we cannot make causal claims about the personality--mental health association.

Third, Study~3 used the BFI-10 \citep{rammstedt2007measuring}, which has lower internal consistency ($\alpha \approx .65$) than the BFI-44 used in Studies~1 and~2. Attenuation due to measurement error would bias results \textit{against} personality, and our disattenuation analyses confirmed that reliability-corrected personality effects are, if anything, larger than observed. Nevertheless, the use of an abbreviated instrument in the largest study introduces measurement imprecision that could affect fine-grained comparisons.

Fourth, we analyzed secondary data not originally designed for this specific comparison. Sensing modalities differ across studies---Studies~1 and~3 relied on smartphone sensing, while Study~2 included Fitbit wearable and communication metadata---making direct comparisons of sensing ``dosage'' across datasets imperfect. The absence of GPS and social media features in some datasets means we cannot rule out the possibility that unmeasured behavioral domains carry stronger signal.

Fifth, all mental health outcomes are self-report symptom scales, not clinical diagnoses. Questionnaire cutoffs (e.g., CES-D $\geq$ 16, BDI-II $\geq$ 14) approximate but do not equate to DSM-5 diagnostic criteria, and prediction performance in clinical diagnostic contexts may differ.

Finally, Study~1's small sample ($N = 27$-28) was treated as a discovery dataset throughout. Its results are informative for hypothesis generation but should not be interpreted as confirmatory evidence.

% ══════════════════════════════════════════════════════════════════════
\section{Conclusion}

Across three universities and 1,559 participants, a five-minute personality questionnaire consistently and substantially outperformed weeks of continuous passive behavioral sensing for predicting depression, anxiety, stress, and loneliness. Personality questionnaires explained meaningful variance in mental health outcomes ($R^2 = .09$--.52), whereas sensing-only models performed at or below chance in 12 of 15 comparisons. Sensing added statistically significant incremental validity in only 1 of 8 comparisons after false discovery rate correction, and the effect was clinically negligible ($\Delta R^2 = .024$). More complex machine learning models (multilayer perceptrons, random forests with Optuna tuning) could not compensate for the limited between-person signal in behavioral features. These findings suggest that the field's enthusiasm for digital phenotyping as a screening tool has outpaced its empirical foundations, and that passive sensing's value may lie not in replacing questionnaires for population-level prediction, but in complementing them through personalized, within-person monitoring of high-risk individuals identified by brief self-report. We recommend that future sensing studies include a personality questionnaire baseline as standard practice, providing the simple benchmark against which the incremental contribution of continuous behavioral data can be honestly evaluated.

% ══════════════════════════════════════════════════════════════════════
% END MATTER
% ══════════════════════════════════════════════════════════════════════

\section*{Data Availability}

Study~1 (StudentLife) data are available from the StudentLife project (\url{https://studentlife.cs.dartmouth.edu/}). Study~2 (NetHealth) data are available via \url{https://nethealth.nd.edu/} with registration. Study~3 (GLOBEM) data are available via PhysioNet (\url{https://physionet.org/content/globem/}) with credentialed access and a data use agreement. Analysis code is available at [GitHub URL].

\section*{Ethics Statement}

All three studies received institutional ethical approval from their respective universities (Dartmouth College, University of Notre Dame, University of Washington). Informed consent was obtained from all participants in the original studies. This secondary analysis was conducted under [NUS IRB reference].

\section*{Author Contributions (CRediT)}

\textbf{[Author A]:} Conceptualization, Data curation, Formal analysis, Methodology, Software, Visualization, Writing -- original draft. \textbf{Cyrus S.~H.~Ho:} Supervision, Writing -- review \& editing.

\section*{Conflicts of Interest}

The authors declare no competing interests.

\section*{Funding}

[To be added]

% ══════════════════════════════════════════════════════════════════════
\bibliographystyle{apalike}
\bibliography{references}

% ══════════════════════════════════════════════════════════════════════
% TABLES & FIGURES (at end, or inline — depending on journal)
% Table 1: Three-study dataset overview
% Table 2: Personality × MH correlations (3 studies)
% Table 3: ML prediction R^2
% Table 4: Meta-analysis results
% Table 5: Clinical classification AUC
% Table 6: Incremental validity F-tests
% Table 7: SHAP vs traditional rankings
% Table 8: Demographic controls
% Figure 1: Forest plots
% Figure 2: AUC comparison
% Figure 3: SHAP ranking agreement

\end{document}
